% SPDX-License-Identifier: CC-BY-4.0
% Copyright (c) 2026 Martin Schr\"oder <info@swedishembedded.com>

\section{Replication on eight authors}\label{sec:replication}

The study of \cref{sec:results} has three authors, one training run each and an examination
of \ExamItems\ items. Its main comparison, the adapter against retrieval, was repeated on
eight new authors under a protocol written and frozen before any of their corpora existed.
This section gives that protocol; its results follow it.

\subsection{Protocol}\label{sec:protocol}

\paragraph{Authors and corpora.} Eight authors were drawn by seed, each with its own beliefs,
lexicon and situations, by the procedure of \cref{sec:benchmark}. Each corpus plans 60 essays
of three episodes each and at least 700 words; a writer, a reader and a miner model
(\cref{sec:reproducibility}) write the essays, read the decisions in them and mine their
briefings. The protocol fixes these models, the number of
attempts per essay (3) and a ceiling on requests (600) and output tokens per author.

\paragraph{Examination.} Each author's examination has 480 items: 120 in each of the four
cells of withheld and seen topics and combinations of \cref{sec:benchmark}. The number was
chosen before the corpora existed, from the first study's author 101: of its 160 items, 37
were discordant between the adapter and retrieval, 23 of them in the adapter's favour. At that
rate of discordance (0.23) and share of wins (0.62), 480 items give an exact sign test at
$\alpha = 0.05$ a power of 0.80 for one author.

\paragraph{Arms.} The arms of \cref{sec:arms} are asked of every author, with three more
retrieval arms: two and twelve paragraphs instead of six, and retrieval of two, six or twelve
whole decision records instead of paragraphs. The adapter is trained three times per author
with seeds the protocol draws, with the hyperparameters of \cref{sec:lora} except the base's precision (below); an author's
adapter accuracy is the mean over the three. Two ablations train once each on a part of the
briefings: the reader's own wording only, and the fixed wording only.

\paragraph{Comparisons.} The primary comparison is the adapter against retrieval of six
paragraphs. The confirmatory family has ten comparisons of the adapter: against
retrieval of six paragraphs and of six records, base, named, principles, principles with
verdicts, the two ablations, and the trivial baselines that always take the first-listed or
the majority option; each is a sign test over authors (an author counts
as better when its paired gain is positive), with Holm's correction over the family at
$\alpha = 0.05$. The other comparisons are exploratory and reported without correction. Each
author's gain is the paired difference in accuracy, with an interval clustered by item.

\paragraph{Decoding.} Greedy decoding is primary. As a secondary decoding, every item is
asked eight more times at temperature 0.7 with a seed per draw.

\paragraph{Amendment.} The protocol held the base model in fp32 during training, which for an
8B model is 32~GB of weights and does not fit one 24~GB card. Before any adapter was trained,
the protocol was amended to hold the training base in int8 with one scale per 32 weights. The
amendment is recorded in the protocol with the digest of the version it replaces; corpora made
under the earlier digest are accepted because the change does not touch them.

\paragraph{Serving.} Every arm is asked through one serving path: a local server holds the
base, or the base with one adapter, with int8 weights and an int8 key--value cache, and
decodes 128 sequences together. The examination keeps all 128 decoding slots busy. Batched
decoding is not bit-identical to decoding one sequence at a time; on 532 items of two authors'
base arms, the batched server and the sequential in-process model chose the same option on 432
(81\%). Because all arms share the batched path, the comparisons between arms are unaffected
by this difference, but absolute accuracies are not comparable with \cref{sec:results} to
better than a few points.

\IfFileExists{study_macros.tex}{% SPDX-License-Identifier: CC-BY-4.0
% Copyright (c) 2026 Martin Schr\"oder <info@swedishembedded.com>

% Read only when the study's report has been built into study_macros.tex and
% the study_*.tex tables (scripts/study_tables.py); every number is theirs.
\input{study_macros}

\subsection{Results}\label{sec:replication-results}

\paragraph{Primary comparison.} The adapter was more accurate than retrieval of six
paragraphs for \StudyPrimaryImproved\ of \StudyPrimaryMeasured\ authors (sign test over
authors $p = \StudyPrimaryP$, Holm-corrected $p = \StudyPrimaryPHolm$); its mean paired gain
was \StudyPrimaryMeanGain\ points. Of the \StudyConfirmatory\ confirmatory comparisons,
\StudyConfirmed\ were confirmed at $\alpha = 0.05$ after Holm's correction
(\cref{tab:study-comparisons}). Each author's gain, with its interval clustered by item, is
in \cref{tab:study-authors}.

\begin{table}[t]
\centering\small
\input{study_comparisons}
\caption{Every comparison the protocol names, over the \StudyAuthors\ authors. An author is
better when its paired gain is positive. The confirmatory family comes first, with Holm's
correction; the exploratory comparisons below it are uncorrected.}
\label{tab:study-comparisons}
\end{table}

\begin{table}[t]
\centering\small
\input{study_authors}
\caption{Each author's accuracy under retrieval and the adapter (mean over its
\StudyTrainingSeeds\ training seeds), on \StudyItemsPerAuthor\ items, with the paired gain
and its 95\% interval clustered by item, the items the adapter alone answered right and
wrong, and the exact sign test over items.}
\label{tab:study-authors}
\end{table}

\paragraph{All arms.} \Cref{tab:study-arms} gives each arm's accuracy averaged over the
authors, with the lowest and highest author. \Cref{tab:study-spread} shows how much an
author's adapter varies with the training seed, and \cref{tab:study-corpora} what each
corpus held: \StudyArticlesKept\ essays and \StudyDecisionsKept\ decisions in all.

\begin{table}[t]
\centering\small
\input{study_arms}
\caption{Accuracy of every arm, averaged over the authors, with the lowest and highest
author.}
\label{tab:study-arms}
\end{table}

\begin{table}[t]
\centering\small
\input{study_spread}
\caption{The adapter's accuracy across its training seeds, per author.}
\label{tab:study-spread}
\end{table}

\begin{table}[t]
\centering\small
\input{study_corpora}
\caption{Each author's corpus: essays planned and kept, cases read, decisions kept and the
requests the writer, reader and miner made.}
\label{tab:study-corpora}
\end{table}
}{}
